\documentclass[10pt]{article}
\usepackage[margin=20mm]{geometry}
\usepackage{amsmath,amssymb,fontspec,unicode-math,xurl,hyperref}
\IfFontExistsTF{DejaVu Serif}{\setmainfont{DejaVu Serif}}{\setmainfont{Latin Modern Roman}}
\setmathfont{latinmodern-math.otf}
\setlength{\parindent}{0pt}\setlength{\parskip}{6pt}\sloppy
\begin{document}
Scientific Verification Report | I25-40104\par
Position: Solve Layerwise Linear Models First to Understand Neural Dynamical Phenomena (Neural Collapse, Emergence, Lazy/Rich Regime, and Grokking)\par
\section*{1. Identity, version and outcome}
Nam et al., ICML 2025 position track, PMLR 267:81897-81929. Primary source: the 33-page PMLR proceedings PDF downloaded from the publisher-linked mlresearch repository [S1]. Page references in this report are physical PDF pages 1-33, which match the printed page numbers; they are not proceedings-wide page numbers. SHA-256: efee741d64336879fd3811427574121fe766889f5478000455c27111035db2b2. The arXiv record lists v2, 26 May 2025, as the latest revision retrieved; targeted comparison of Appendix H found the same formulas [S2]. Byte identity with the arXiv PDF is not asserted.\par
Outcome: Major for Problem 1 (F01); Problem 2 (F02) remains Minor. Theorem 3’s exact-time solution is an expressly advertised part of the Section 7 toy-model contribution and violates the paper’s stated gradient flow. This is a core-related invalid derivation under the clarified grading standard. The released Figure 6(b) axis finding is unchanged. The existing calculations and evidence are reused; the precise surviving endpoint conclusions and original-domain qualifications are stated below. Human scientific decision and review remain pending; author adoption has not been requested and publication is disabled.\par
\section*{2. Scope, definitions and checked reasoning}
The review covered the main text (pp.1-9), relevant definitions and appendices A-I (pp.16-33), and the released Figure 6/8 implementation. Bibliography pp.10-15 was used for locating sources, not exhaustively audited. The source code calls the current Figure 8 experiment Figure 11. Source inspection and analytic checks were performed; no author code, model training or empirical reproduction was run.\par
Material | Actual check | Outcome\par
Section 3; Appendix C.1-C.3, pp.3,19-20 | Differentiate the MSE; verify diagonal and matrix Gram-balance invariants. | No issue found in the checked equations.\par
Section 4; Appendices C.2 and D, pp.3-4,20,22-23 | Check logistic mode solution and mutually exclusive skill decomposition. | Supported for positive balanced modes and the declared feature-support model.\par
Section 5; C.4 and E, pp.5,20-21,24 | Check aligned mode calculation; compare informal NC discussion with NC1-NC4 definitions. | Scope clarification: exact task alignment and balanced classes matter; rank alone is not an ETF proof.\par
Section 6; Appendix F, pp.6,25-27 | Check imbalance-mode intuition, endpoint NTK metric and code mappings. | Theory checks supported within scope; released axis issue below.\par
Section 7; G-H, pp.7-8,28-32 | Reconstruct scalar gradient flow, inspect five MNIST configurations and general-initialization formula. | Theorem 3 time-scale defect; other experiment settings match the inspected code.\par
Sections 8-10; Appendix I, pp.9,33 | Check scope and downstream dependence. | Paper expressly does not claim that greedy dynamics guarantees generalization.\par
For the scalar toy model, let p be a positive integer, S>0 the target coefficient, Z>0 the normalization, and a\_i,b\_i finite real parameters. Let theta=(sum\_i a\_i b\_i)/Z, S0=theta(0), Sigma0=sum\_i(a\_i(0)\textasciicircum{}2+b\_i(0)\textasciicircum{}2)/(2Z), A=Sigma0+S0, B=Sigma0-S0, and D=sqrt(S\textasciicircum{}2+AB). The positive-Z domain is necessary for the norm inequalities used by the paper. Standard normal scalar input makes the expected loss agree exactly with its Eq.(71).\par
The diagonal equations give d(a\_i\textasciicircum{}2-b\_i\textasciicircum{}2)/dt=0 and, for a\_i=b\_i, d(a\_i b\_i)/dt=2 E[x\_i\textasciicircum{}2] a\_i b\_i(S\_i-a\_i b\_i). This verifies the positive-target logistic calculation. The exact matrix decoupling in C.4 uses task-aligned balanced initialization; small arbitrary initialization alone is not that assumption. The neural-collapse passage is explicitly informal. Its low-rank intuition was not treated as a proof of every NC condition. The cosine NTK distance in F.2 measures directional change: it is zero when K\_t=c K\_0 for c>0. This acknowledged endpoint diagnostic is not a complete test of frozen-kernel dynamics.\par
\section*{3. Problem 1: Theorem 3 doubles the stated gradient-flow speed}
\subsection*{Source claim and suspect step}
Theorem 3, p.30, Eqs.(72)-(74), claims an exact solution of the loss in Eq.(71) and the gradient-flow vector field in Eq.(77). Eq.(74), repeated as Eq.(96) on p.32, contains exp(-4Dt/Z). The conversion to the Riccati equation in Eq.(89), p.31, swaps two coefficients in its first bracket and doubles the coefficient of its final quadratic. This reaches the stated theorem, so the issue affects an exact mathematical result rather than only typography. The actual PDF equations were visually checked [S1]; arXiv v2 retains them [S2].\par
\[\mathcal L=\tfrac12(\theta-S)^2,\quad
\dot a_i=\frac{S-\theta}{Z}b_i,\quad
\dot b_i=\frac{S-\theta}{Z}a_i.\]\par
\subsection*{Complete analytic check}
The paper’s definition gamma=exp[(2/Z) integral\_0\textasciicircum{}t (S-theta(s)) ds] gives its parameter representation correctly. Set u\_i=(a\_i(0)+b\_i(0))/2 and v\_i=(a\_i(0)-b\_i(0))/2. Then a\_i=u\_i sqrt(gamma)+v\_i/sqrt(gamma) and b\_i=u\_i sqrt(gamma)-v\_i/sqrt(gamma). Multiplication and summation produce the following identity and differential equation:\par
\[\theta=\frac{A\gamma-B\gamma^{-1}}{2},\qquad
\dot\gamma=-\frac{1}{Z}\left(A\gamma^2-2S\gamma-B\right).\] Equivalently, the first bracket of corrected Eq.(89) is S0(gamma+1/gamma)/2 + Sigma0(gamma-1/gamma)/2 - S, multiplied by -2gamma/Z. For A>0, let r\_minus=(S-D)/A and r\_plus=(S+D)/A. Separating the Riccati equation yields:\par
\[\frac{\gamma(t)-r_+}{\gamma(t)-r_-}
 =\frac{1-r_+}{1-r_-}e^{-2Dt/Z},
\qquad
\gamma(t)=\frac{r_+-q(t)r_-}{1-q(t)},\quad
q(t)=\frac{1-r_+}{1-r_-}e^{-2Dt/Z}.\] This expression is algebraically the published rational formula with each 4 in its exponential rate replaced by 2. It starts at gamma=1. Its derivative is exactly the Riccati equation above, and substitution in the parameter representation recovers both original gradient equations. Since A>0, B>=0 and D>=S>0, gamma remains positive and converges monotonically to r\_plus. The finite-equilibrium formula survives unchanged on this domain.\par
\subsection*{Fully specified counterexample}
Take p=1, Z=1, S=1, a(0)=1 and b(0)=0, with standard normal input. All setup assumptions hold, S0=0 and Sigma0=1/2>|S0|, so this is not a singular boundary case. The original gradient field gives a’(0)=0 and b’(0)=1. Differentiating the published formula gives a’(0)=0 and b’(0)=2. Independent SymPy 1.13.3 calculation confirms those derivatives and a nonzero original-ODE residual; after the correction the residual simplifies identically to zero [S4].\par
Quantity | Published time formula | Corrected original gradient flow\par
a’(0), b’(0) | 0, 2 | 0, 1\par
gamma-ODE residual at t=0 | 2 | 0\par
a(1), b(1) | 1.25993486863, 0.766443652972 | 1.17953536267, 0.625542701811\par
theta(1) | 0.965669083217 | 0.737849737647\par
An elementary cross-check takes p=Z=1, S=2 and a(0)=b(0)=1. The correct solution is a(t)=b(t)=sqrt[2/(1+exp(-4t))], whereas the theorem gives sqrt[2/(1+exp(-8t))]. Their initial derivatives are 1 and 2, respectively.\par
\subsection*{Alternative reading, impact and verified repair domain}
Judgment: Major (calibrated from Minor). Section 1.1, p.2, explicitly lists solving the Section 7 toy model as a contribution. Section 7.1, pp.7-8, advertises exact dynamics for general initialization and invokes Theorem 3. The dependency is the toy-model contribution, its exact general-initialization trajectory, and then its limiting amplification factor gamma\_plus and weight-to-target interpretation. The existing counterexample invalidates the exact-time trajectory node, so the derivation is closely tied to a core contribution.\par
A doubled learning rate or loss without the factor 1/2 would give a coherent different clock, but that convention is not consistently used in this paper: Eqs.(2), (6), (9), (16)-(17), and (71)/(77) fix the original normalization. On the verified A>0, S>0, Z>0 domain, w\_published(t)=w\_correct(2t). Thus values at a stated time are values at twice that time, not generally twice the correct values; derivatives are twice the correct derivatives evaluated at 2t, and attained threshold times are halved. The parameter-space path, endpoint, and ratios of threshold times subject to this same rescaling survive. Stationary trajectories are unaffected.\par
The strongest counterevidence is the actual downstream use on p.8: the general-initialization rich/lazy argument uses gamma\_plus and limiting products, not the erroneous exponential. Those endpoint expressions, the weight-to-target criterion and Remark 5 survive on A>0, as do the separately correct special-case equations on p.7. No evidence ties the erroneous formula to the empirical Figure 8 curves. These limits confine the demonstrated damage to exact-time predictions; they do not make the advertised exact-dynamics derivation peripheral. The Major judgment does not depend on invalidating the broader position or observed scores.\par
Scope clarification for ’any initialization’: the finite equilibrium expression divides by A. For A=0, all initial pairs satisfy a\_i(0)=-b\_i(0). The correct separate branch is gamma(t)=(1+Sigma0/S)exp(2St/Z)-Sigma0/S. All parameters then approach zero, theta approaches zero and the loss approaches S\textasciicircum{}2/2. For p=S=Z=1 and initial pair (1,-1), gamma=2exp(2t)-1, a=1/sqrt(gamma), b=-a. Both gradient residuals simplify to zero [S4]. Thus the verified finite-equilibrium repair explicitly assumes A>0, with this boundary treated separately; it does not assert convergence to the target for every initialization.\par
\section*{4. Problem 2: released Figure 6(b) axis disagrees with its parameter grid}
\subsection*{Source claim and complete static check}
Figure 6(b), p.6, labels its vertical axis by sigma1-sigma2 and shows +5 at the top and -5 at the bottom; the PDF rendering was inspected. Appendix F.4, p.27, explains initialization using these standard deviations. At pinned code commit 95c2c4a4609b386c3f8c1af5504d3ac549b82620, fig6b\_run.py line 18 instead sweeps lambda from -0.5 to +0.5. Line 40 calls get\_lambda\_unbalanced with scale=0.1; utils.py lines 327-332 set sigma1=0.1+max(lambda,0) and sigma2=0.1+max(-lambda,0). Hence their difference equals lambda exactly [S3].\par
The run script sorts rows in increasing lambda order (lines 62-64), fills that order (lines 98-100), and saves it without reversal (lines 144-145). The separate plot script uses labels [5,0,-5] at rows [0,rows//2,rows-1] (lines 21-22) and explicitly names sigma1-sigma2 (line 42). There is no intervening sign reversal or factor ten scaling.\par
Saved row | Actual initialization SDs | Actual difference | Displayed label\par
First: lambda=-0.5 | sigma1=0.1; sigma2=0.6 | -0.5 | +5\par
Last: lambda=+0.5 | sigma1=0.6; sigma2=0.1 | +0.5 | -5\par
The symmetrized initialization itself is correct: W2 W1 = sigma2[B,-B] sigma1[A;A] = 0. The defect is the quantitative axis mapping, which reverses the identity of the heavier layer and exaggerates the magnitude tenfold.\par
\subsection*{Counterevidence, impact and local repair}
The current repository contains no original heatmap arrays or run logs binding this checkout to the published figure. A different original pipeline remains possible. Therefore the released-code discrepancy is demonstrated; the exact provenance and quantitative correction of the published image remain Needs Evidence. This distinction prevents a code observation from becoming an unsupported claim about which experiments were run.\par
For the saved row order, derive all tick labels from the actual lambda grid: the endpoints must be -0.5 and +0.5. Reversing the matrix would instead permit +0.5 at the top and -0.5 at the bottom. With 40 grid points there is no exact zero row, so even the center label should be generated from its coordinate or placed by interpolation. This analytically repairs the axis mapping; it has not been validated on original experimental arrays. The code finding is Minor and does not disprove the qualitative existence of rich and lazy regions.\par
\section*{5. Dependence, additional scope checks and remaining work}
The MNIST code’s four-layer width-512 tanh architecture, MSE, 1000 training images, Adam settings and five scaling configurations agree with Appendix G.3 (pp.28-29). Its test routine sets training mode, but the model has neither dropout nor batch normalization, so this alone is not a prediction defect. Seeds and repeated-run statistics are not supplied. Reproduction would be needed to assess curve variability and effect robustness; none was performed.\par
A secondary Figure 6(a) calibration question was retained in the coverage ledger rather than promoted to a separate published-result defect. The active get\_lambda\_balanced implementation makes the two product singular values sqrt(s), giving Frobenius norm sqrt(2s), while fig6a\_run.py also scales targets by 1/s. Consequently the plotted nominal parameter s cannot simply be read as the product Frobenius norm; the intended ’weight-target ratio’ and original figure pipeline require clarification. This pending question does not increase the supported severity.\par
The arXiv v2 record was inspected and the Appendix H error remains in its HTML. The public repository’s current tip, eight-commit history metadata, issues, tags and releases were queried; no correction was identified in those records. A full historical diff and an exhaustive search for every external erratum were not performed. No existing correction was attributed to this review, and no claim of priority for the observations is made.\par
Item | Issue judgment | Repair verification | Adoption / remaining work\par
Problem 1 | Major: invalid exact-time dynamics in the stated toy-model contribution | General analytic reconstruction; symbolic counterexample and boundary checks | Not requested. Human decision pending.\par
A=0 boundary | Scope clarification | Separate solution verified analytically and symbolically | State domain in any corrected theorem.\par
Problem 2 | Minor: released plotting axis | Mapping derived exactly; no empirical rerender | Original arrays and figure-generation provenance needed.\par
Other covered claims | No issue found within specified scope | Checks listed in Section 2 and coverage.json | No whole-paper correctness claim.\par
\section*{6. Consulted references and stable source identifiers}
[S1] Nam, Y., Lee, S. H., Domine, C. C. J., Park, Y., London, C., Choi, W., Goring, N. A., and Lee, S. Position: Solve Layerwise Linear Models First to Understand Neural Dynamical Phenomena. ICML 2025, PMLR 267:81897-81929. Publisher record: https://proceedings.mlr.press/v267/nam25a.html . Exact downloaded PDF: https://raw.githubusercontent.com/mlresearch/v267/main/assets/nam25a/nam25a.pdf .\par
Same authors and title. arXiv:2502.21009v2, revised 26 May 2025. Record: https://arxiv.org/abs/2502.21009 ; version-pinned comparison: https://arxiv.org/html/2502.21009v2 . Targeted formula comparison only.\par
yoonsoonam119/linear\_first, commit 95c2c4a4609b386c3f8c1af5504d3ac549b82620, dated 23 May 2025. https://github.com/yoonsoonam119/linear\_first . Pinned raw files, file hashes and history-query snapshots accompany the report under sources/0006.\par
Independent reviewer calculations, sources/0006/independent\_checks.py and independent\_checks.json. SymPy 1.13.3. These are reviewer-authored symbolic checks, not author-code execution or empirical reproduction.\par
\section*{Supporting checks}
Executed independent check (STAGING expands to the assigned root):
[local archive path omitted] -B -X utf8
STAGING\textbackslash{}sources\textbackslash{}0006\textbackslash{}independent\_checks.py\par
Observed: exponent=4 -> b\_dot0=2; ODE residual at0=2.
Observed: exponent=2 -> b\_dot0=1; residual identically zero.
Anti-aligned branch: both ODE residuals=0; loss limit=1/2.\par
The full executed source and exact output are preserved as S4. Source extraction used the installed pdftotext; PDF page count was confirmed with pdfinfo. Embedded form-feed characters in a figure cause naive text splitting to count 35 chunks, while pdfinfo and printed page numbers correctly give 33 pages.\par
OpenReview presented a verification wall, so its discussion and attachment history were not audited. Pinned GitHub web views were unavailable in one route; raw-source retrieval succeeded. An initial Python capability probe lost its string quotes and raised NameError; a here-string probe succeeded. No failed mutation was blindly retried.\par
No independent human review or multi-model benchmark is claimed. No human scientific approval, author communication or adoption is asserted.\par
\end{document}
